\par\noindent\textbf{(1)} 临界条件 \((\partial_VP)_T=-RT/(V-b)^2+2a/V^3=0\) 与 \((\partial_V^2P)_T=2RT/(V-b)^3-6a/V^4=0\) 联立给 \(V_c=3b\)，\(T_c=8a/(27Rb)\)，\(P_c=a/(27b^2)\)。必须同时满足两个导数为零，非任意拐点。
\par\noindent\textbf{(2)} 令 \(p=P/P_c,t=T/T_c,v=V/V_c\)，代入得到 \(p=8t/(3v-1)-3/v^2\)。临界压缩因子 \(Z_c=P_cV_c/(RT_c)=3/8\)。它是这一平均场模型的值，并非所有实际流体严格通用。
\par\noindent\textbf{(3)} 在 \(t=1,v=1+x\) 附近展开 \(8/(2+3x)-3/(1+x)^2=1-\tfrac32x^3+O(x^4)\)。因为线性、二次项消失，\(|P-P_c|\propto|V-V_c|^3\)，所以平均场临界指数 \(\delta=3\)。
\par\noindent\textbf{(4)} 等温压缩率 \(\kappa_T=-(1/V)(\partial V/\partial P)_T\)。在临界等容线 \(v=1,t=1+\tau\) 上，\((\partial p/\partial v)_t=-6\tau\)，故 \(\kappa_T\propto|\tau|^{-1}\)，平均场 \(\gamma=1\)。接近真实临界点时涨落改变指数，van der Waals并非精确临界理论。
